\chapter{Introduction}
\label{cha:intro}

\section{Forecasting a plant you have never seen}
\label{sec:motivation}

A photovoltaic installation converts a stochastic, weather-driven input into an
electrical output that a grid operator has to schedule around. The forecasting problem
that follows is a short-horizon regression: given the recent power history of a plant and
whatever exogenous information is available, predict the next few hours of generation.
Framed that way it looks like a solved family of problems, and for a \emph{single,
instrumented, long-lived} plant it very nearly is --- a gradient-boosted model over lagged
power and weather covariates is a strong forecaster, and a supervised transformer trained
on years of that plant's own history is stronger still.

The formulation breaks down at the point where it matters operationally. Photovoltaic
capacity is added continuously and in small increments: rooftop systems, community
installations, new inverters on an existing site. A plant that was energised last month has
no multi-year history to fit a model to. If every installation requires its own training
run, the forecasting stack scales linearly in the number of assets and degrades exactly
where the fleet is growing fastest. The deployment-realistic question is therefore not
``how well can we forecast a plant we have studied'' but

\begin{quote}
  \emph{how well can we forecast a plant we have never seen, given two weeks of its history
  and nothing else?}
\end{quote}

This is a generalisation question, not a photovoltaic-engineering question, and it is the
question this thesis is built around. The distinction matters for how the work is
evaluated. Under the usual protocol a model is trained and tested on the same plants,
separated only in time, and reported accuracy conflates two very different abilities:
learning the idiosyncrasies of a specific installation, and learning something transferable
about how irradiance, cloud cover and time of day map onto power. Splitting by \emph{plant
identity} rather than by time isolates the second ability, and it is the only split under
which the deployment question can be answered honestly. Throughout this thesis, train,
validation and test sets are disjoint sets of \emph{plants}; no test plant contributes a
single row, a single frame or a single normalisation statistic to anything that is fitted.

\subsection{What makes the signal hard}

It is worth being precise about the difficulty, because ``predict solar power'' sounds easier
than it is. The signal decomposes into two parts with completely different character.

The first is deterministic: the position of the sun, and therefore the clear-sky irradiance
reaching a fixed panel, is known analytically for all future time. Any model with a calendar
and a latitude reproduces it. It also accounts for the majority of the raw variance, which
means that a forecast evaluated against a reference that does not know it will look excellent
while having learned nothing.

The second is the modulation of that envelope by the atmosphere --- overwhelmingly by cloud.
It is intermittent, spatially structured, advects at wind speed, and changes the output of an
array by most of its range within minutes. It is not well captured by hourly reanalysis
weather, which is what numerical covariates provide: an hourly areal cloud fraction says
little about whether the sun is occluded at one rooftop at 14:20. This is the part a
forecaster is actually paid to predict, and it is the part for which the numerical track is
an impoverished description.

That gap is the opening for a visual modality. Cloud fields are directly observable in
satellite imagery, at a spatial and temporal resolution that reanalysis does not approach,
and they advect predictably enough over a few hours that what is visible upwind now is
overhead later. Chapter~\ref{cha:dataset} quantifies the residual information available:
conditioning on cloud state changes median output by roughly 0.09 in normalised power
\emph{after} conditioning on irradiance. Whether any current architecture converts that
availability into forecast skill is, in the end, the question this thesis is about.

\section{Why foundation models, and why multimodal}
\label{sec:why-fm}

Two developments make the cross-plant question tractable in a way it was not a few years
ago.

The first is the arrival of \emph{time-series foundation models}: transformers pretrained
on large, heterogeneous corpora of time series that forecast unseen series zero-shot, with
no per-series fitting~\cite{chronos2,timesfm,tirex}. Their promise is precisely the
promise this problem needs --- transfer to a series that was not in the training set. The
second is the maturation of \emph{video foundation models} trained by self-supervised
prediction rather than reconstruction~\cite{vjepa2}. This matters because the dominant
source of short-horizon photovoltaic uncertainty is cloud advection, a spatiotemporal
process that is visible in satellite imagery hours before it registers in the power signal.
A model that reads only the power history is, in the ramp regime, structurally blind: the
information that would let it anticipate a cloud edge exists, but not in the modality it
consumes.

Combining the two is not novel in itself; a growing literature attaches visual encoders to
time-series forecasters~\cite{timevlm,unicast,solarvlm,crossvivit}. What is contested is
\emph{how deeply} the two modalities should be coupled. The prevailing designs are shallow:
a visual encoder produces a summary vector which is concatenated to the series
representation (late fusion), or injected as a soft prompt~\cite{unicast}, or used to
condition an adapter on an otherwise frozen backbone~\cite{cora}. In each case the two
modalities meet once, after each has been processed independently. The alternative
investigated here is to let them meet \emph{inside} the temporal operator, as neighbouring
tokens in one sequence, so that the model computes cross-modal interactions at every layer
rather than a single cross-modal readout at the end.

The claim under test is therefore:

\begin{quote}
  A frozen multimodal foundation model stack (Chronos-2 and V-JEPA~2.1) with deep
  token-level fusion achieves cross-plant photovoltaic power forecasting on disjoint test
  plants, beating late fusion, unimodal foundation models, and domain-specific
  architectures.
\end{quote}

Every baseline in Chapter~\ref{cha:baselines} exists to falsify one part of that sentence.
The design principle behind the benchmark is that a reader must not be able to say ``the
gain could come from $X$, and $X$ was never tested''.

\section{Research question and hypotheses}
\label{sec:rq}

The research question, stated in the terms the thesis is evaluated in:

\begin{quote}
  Can a frozen multimodal foundation model stack (a time-series foundation model and a
  vision foundation model) achieve cross-plant photovoltaic power forecasting on disjoint
  test plants by deep token-level fusion, rather than by late fusion or by domain-specific
  architectures?
\end{quote}

It decomposes into a ladder of hypotheses, each of which is answered by a specific
experiment, and each of which is falsifiable independently of the ones above it.

\begin{itemize}
  \item \textbf{H0 --- Foundation-model anchor.} A time-series foundation model applied
    zero-shot is a meaningful reference point for this task. Tested in
    Chapter~\ref{cha:baselines} against classical and supervised alternatives.
  \item \textbf{H1 --- Vision helps.} Adding a visual stream through a late-fusion adapter
    improves over an otherwise identical time-series-only model. Tested by the modality
    ablation of Chapter~\ref{cha:expdesign}.
  \item \textbf{H2 --- Fusion depth matters.} Selective temporal interleaving of visual
    tokens improves over late fusion, with the same backbones and the same data. This is
    the architectural contribution of Chapter~\ref{cha:method}.
  \item \textbf{H3 --- Cross-plant transfer is achievable.} A model can forecast completely
    disjoint plants from a short history, relying on structure learned from other plants.
    This is the evaluation protocol of Chapter~\ref{cha:protocol}, not an ablation.
  \item \textbf{H4 --- Retrieval is a competitive alternative.} Retrieval augmentation on a
    frozen backbone closes the gap to full fine-tuning. Tested in
    Chapter~\ref{cha:baselines}.
\end{itemize}

\section{Contributions}
\label{sec:contributions}

\paragraph{A dataset of record for multimodal cross-plant photovoltaic forecasting.}
Four public sources --- UK domestic generation~\cite{ukpv}, EUMETSAT SEVIRI high-resolution
visible imagery~\cite{eumetsat}, NREL PVDAQ systems with GOES-16 crops~\cite{pvdaq}, and
Open-Meteo reanalysis covariates~\cite{openmeteo} --- are fused into one flat table of
$1{,}337{,}654$ rows over 110 sites, plus a packed image archive addressed by a single
timestamp-exact pointer. The construction, the quality flags and the exact plant partition
are given in Chapter~\ref{cha:dataset}.

\paragraph{A 28-model benchmark under one fairness contract.}
All models consume the same tensor contract, the same 14-day history and 6-hour horizon
defined in \emph{physical} time, the same plant-disjoint split, and the same
capacity-normalised metrics. The suite spans seven tiers, from naive persistence to
photovoltaic-specialised multimodal networks. To the best of this work's knowledge it is
the largest single-protocol comparison of time-series foundation models,
frozen-backbone adaptation methods and multimodal forecasters on photovoltaic data.

\paragraph{MMTSFM, an architecture for deep multimodal fusion on frozen backbones.}
Two mechanisms are contributed. \emph{Selective temporal interleaving} inserts visual
summary tokens into the time-series sequence only inside a short refinement window, leaving
the macro-context purely numerical; the token overhead is proportional to the number of
visual tokens rather than to context length, on the order of two percent. \emph{Causal
Grassmann mixing} replaces $O(L^2)$ temporal attention with an $O(L)$ operator that encodes
the transition between hidden states as a two-dimensional subspace via a Pl\"ucker embedding,
aggregated over multiple temporal offsets. A matched diagnostic twin using ordinary causal
self-attention is retained so that any measured difference is attributable to the inductive
bias and not to the training schedule.

\paragraph{A diagnosis of where the headroom actually is.}
Two measurements from the benchmark reorient the problem. The gap between a
privileged-information oracle variant of Chronos-2 and its honest fine-tuned counterpart is
$0.173$ skill --- larger than any architectural difference observed in the entire suite,
and evidence that \emph{conditioning}, not architecture, is the binding constraint.
Separately, the strongest multimodal model in the suite uses images synthesised from the
series itself, while every model consuming real satellite frames ranks mid-pack; the
distinction between visual \emph{inductive bias} and visual \emph{information} is developed
in Chapters~\ref{cha:baselines} and~\ref{cha:conclusion}.

\section{Scope and non-goals}
\label{sec:scope}

The framing is deliberately that of a machine-learning contribution with photovoltaics as
the testbed, not a photovoltaic-engineering contribution. Several things are consequently
out of scope, and are listed here so the reader does not go looking for them.

\begin{itemize}
  \item \textbf{No domain physics inside models.} Clear-sky-index conversions and
    irradiance-to-power physical models are excluded from every learned model. The single
    exception is smart persistence, which \emph{is} the physics reference and serves as the
    denominator of the skill score.
  \item \textbf{No few-shot in-context adaptation curves.} Accuracy as a function of $K$
    support days per test plant is a different protocol --- context matching rather than
    zero-shot transfer --- and was excluded by an explicit design decision. The evaluation
    harness supports it as an appendix experiment should a reader require it.
  \item \textbf{No grid-operations or electricity-market analysis.} The output of this work
    is a forecast and its error, not a dispatch decision or an economic valuation.
  \item \textbf{No pre-2025 methods as primary comparators.} Older architectures appear
    only where they are still the community reference (persistence, gradient boosting,
    DLinear as the simplicity check).
  \item \textbf{No dataset construction beyond what is required.} The consolidated dataset
    is treated as read-only once built; no result in this thesis depends on re-processing
    it.
\end{itemize}

\section{A note on what this thesis reports}
\label{sec:reporting-stance}

Two editorial decisions shape how the results are presented, and both are worth declaring at
the outset because they are unusual enough to be mistaken for omissions.

\emph{Negative results are reported as results.} Several of the most informative measurements
in this work are things that did not happen: foundation models did not transfer zero-shot,
real satellite imagery did not outperform synthesised imagery, and pretraining did not beat a
well-matched supervised architecture. Each is stated in the chapter where it was measured,
with the control that establishes it, and none is relegated to a limitations paragraph.

\emph{Provisional numbers are not reported at all.} The experimental campaign for the proposed
architecture was still running at the time of writing. Chapter~\ref{cha:expdesign} therefore
states the arms, the controls, the interpretation rules and the reporting template in full,
and leaves the result cells empty. A number that later moves is worse than no number, because
it propagates silently into every comparison that cites it --- and because a design whose
interpretation is fixed only after the result is known is not a test of anything.

The second decision has a consequence for how the thesis should be read: its empirical weight
sits in Chapters~\ref{cha:dataset} through~\ref{cha:benchmark}, which are complete, rather
than in Chapter~\ref{cha:expdesign}, which is a specification.

\section{Structure of the thesis}
\label{sec:structure}

Chapter~\ref{cha:background} positions the work against the 2025--2026 literature on
time-series foundation models, frozen-backbone adaptation, video foundation models and
multimodal forecasting, and states the gap that motivates the benchmark.
Chapter~\ref{cha:dataset} describes the construction of the dataset of record, its schema,
its quality flags and its plant-disjoint splits. Chapter~\ref{cha:protocol} defines the
evaluation protocol: the fairness rules, the physical-time windows, the metrics and the
scenario battery. Chapter~\ref{cha:baselines} presents the baseline suite, its
implementation, the full leaderboard and the four findings that follow from it.
Chapter~\ref{cha:method} presents MMTSFM in full: both branches, the fusion mechanisms, the
Grassmann operator, the output head and the four-stage training curriculum.
Chapter~\ref{cha:expdesign} gives the experimental design for the model, its ablation and
control battery, and the reporting protocol under which its numbers will be recorded; the
campaign was still running at the time of writing and no provisional figures are stated.
Chapter~\ref{cha:conclusion} answers the research question as far as the current evidence
permits, states the limitations, and lays out the next experiments in priority order.

\par

